%!TEX TS-program = lualatex
%!TEX encoding = UTF-8 Unicode

\documentclass[12pt]{article}
%\usepackage{graphicx}
%	\graphicspath{{/Users/goby/Pictures/teach/153/lab/}} % set of paths to search for images

\usepackage{geometry}
\geometry{letterpaper}                   
\geometry{bottom=1in, left=1.5in}
%\geometry{landscape}                % Activate for for rotated page geometry
\usepackage[parfill]{parskip}    % Activate to begin paragraphs with an empty line rather than an indent
%\usepackage{amssymb}
%\usepackage{mathtools}
%	\everymath{\displaystyle}

%\pagenumbering{gobble}

\usepackage{fontspec}
\setmainfont[Ligatures={Common,TeX}, BoldFont={* Bold}, ItalicFont={* Italic}, Numbers={Proportional, OldStyle}]{Linux Libertine O}
\setsansfont[Scale=MatchLowercase,Ligatures=TeX, Numbers=OldStyle]{Linux Biolinum O}
%\setmonofont[Scale=MatchLowercase]{Inconsolata}
\usepackage{microtype}

% This defines \amper for the fancy ampersand
% to be used in the header. See
% https://tex.stackexchange.com/a/58185/39194
\usepackage{xspace}
\newfontfamily\amperfont[Style=Alternate]{Linux Libertine O}    
\makeatletter
\DeclareRobustCommand{\amper}{{\amperfont\ifx\f@shape\scname\smaller[1.2]\fi\&}\xspace}
\makeatother

\usepackage{unicode-math}
\setmathfont[Scale=MatchLowercase]{Asana-Math.otf}
%\setmathfont{XITS Math}

% To define fonts for particular uses within a document. For example, 
% This sets the Libertine font to use tabular number format for tables.
%\newfontfamily{\tablenumbers}[Numbers={Monospaced}]{Linux Libertine O}
%\newfontfamily{\libertinedisplay}{Linux Libertine Display O}

\usepackage{hanging}

%\usepackage{booktabs}
%\usepackage{multicol}
%\usepackage{tabularx}
%\usepackage{longtable}
%\usepackage{siunitx}
%\usepackage[justification=raggedright, singlelinecheck=off]{caption}
%\captionsetup{labelsep=period} % Removes colon following figure / table number.
%\captionsetup{tablewithin=none}  % Sequential numbering of tables and figures instead of
%\captionsetup{figurewithin=none} % resetting numbers within each chapter (Intro, M&M, etc.)
%\captionsetup[table]{skip=0pt}

\usepackage{array}
\newcolumntype{L}[1]{>{\raggedright\let\newline\\\arraybackslash\hspace{0pt}}p{#1}}
\newcolumntype{C}[1]{>{\centering\let\newline\\\arraybackslash\hspace{0pt}}p{#1}}
\newcolumntype{R}[1]{>{\raggedleft\let\newline\\\arraybackslash\hspace{0pt}}p{#1}}

\usepackage{enumitem}
%\usepackage{hyperref}
%\usepackage{placeins} %PRovides \FloatBarrier to flush all floats before a certain point.
%\usepackage{hanging}
%\usepackage{color}
%\usepackage{calc}

%\usepackage{titling}
%\setlength{\droptitle}{-60pt}
%\posttitle{\par\end{center}}
%\predate{}\postdate{}

\usepackage[sc]{titlesec}


\usepackage{fancyhdr}
\fancyhf{}
\pagestyle{fancy}
\lhead{}
\chead{}
\rhead{\footnotesize pg. \thepage }
\renewcommand{\headrulewidth}{0.4pt}

\fancypagestyle{plain}{%
	\fancyhf{}
	\lhead{\textsc{zo} 678: Ichthyology}
	\rhead{Annotated Bibliography}
	\renewcommand{\headrulewidth}{0pt}
}
	
%\newcommand{\VSpace}{\vspace{0.5\baselineskip}}
%\newcommand{\BigVSpace}{\vspace{2\baselineskip}}

\title{Annotated Bibliography Assignment}
\author{Ichthyology}
\date{}                                           % Activate to display a given date or no date

\begin{document}
%\maketitle
\thispagestyle{plain}


\subsection*{Annotated bibliography (30 points)}

You must write an annotated bibliography of the three papers you have chosen as the basis of your 5-in-5 presentation.

\subsubsection*{Requirements}

\begin{enumerate}[leftmargin=*]

\item Each paper must be properly cited following a standard citation format (see below).  A typical citation format is Author(s) year. Title. Journal name volume number: page range. For example,

\begin{itemize}[label={}, leftmargin=1.5em, itemindent=-1.5em]
\item Karnovsky, N.\,J., K.\,A.~Hobson, S.~Iverson, and G.\,L.~Hunt, Jr. 2008. Seasonal changes in the diets of seabirds in the North Water Polynya: a multiple-indicator approach. Marine Ecology Progress Series 357:291–299.

\end{itemize}

%Andrello, M., D. Mouillot, S. Somot, W. Thuiller, and S. Manel. 2014. Additive effects of climate change on connectivity between marine protected areas and larval supply to fished areas. Biodiversity Research 2014: 1–14.

\item You must type a \emph{thorough} single-spaced two paragraph summary of each paper. %I imagine you should have about 2 citations and summaries per page. 

\item Each summary is worth 10 points. 

\end{enumerate}

\subsubsection*{Other information}

Your summary should focus on the results and discussion as they relate to your topic. Do not simply paraphrase or summarize the abstract. I will grade harshly if I suspect this is the case. Be professional. Read the paper thoroughly and summarize it appropriately. Briefly descibe as part of your second paragraph (1–2 sentences) how you think you will incorporate information from the paper in your presentation. % If you wish, during the first week of the assignment, you may send me one paper with your summary and intent. I will let you know if you are on the right track. You must do this early. I will not do this after 5 pm, one week after I give out this assignment.

\textsc{Due date:} Check the Canvas page for the due date. Upload your completed assignments and the \textsc{pdf} files to the drop box available on the course website. See the syllabus for information on late assignments.

\vspace*{\baselineskip}

Please let me know as soon as possible if you have any questions. If I missed an important detail, I need to let other students know as soon as possible.


\end{document}  